\begin{abstract}
Se implementaron y compararon cuatro algoritmos --- \textit{fuerza bruta}, \textit{programación dinámica} y dos heurísticas \textit{greedy} --- para una variante del problema de la mochila con restricciones de tiempo y energía. Se evaluaron $27$ instancias reproducibles con $n \in \{3,\dots,200\}$, midiendo tiempo de ejecución, memoria del heap y calidad de solución respecto al óptimo. \textit{Fuerza bruta} mostró crecimiento exponencial en tiempo (acotada a $n\leq25$); \textit{programación dinámica} garantizó el óptimo con costo pseudopolinomial y alto uso de memoria; ambas heurísticas \textit{greedy} escalaron polinomialmente con uso de memoria órdenes de magnitud menor. \textit{Greedy ratio} alcanzó consistentemente entre $90\%$ y $100\%$ del óptimo, sin correlación significativa con el tamaño de la instancia; \textit{Greedy max} fue más rápida pero inestable, con calidad entre $\approx 24\%-92\%$ y una dependencia de las variables de entrada.
\end{abstract}